\documentclass[11pt]{article}
\usepackage[T1]{fontenc}
\usepackage{lmodern,microtype}
\usepackage[margin=1in]{geometry}
\usepackage{amsmath,amssymb,amsthm,mathtools,booktabs}
\usepackage[hidelinks]{hyperref}
\usepackage[nameinlink,noabbrev]{cleveref}
\allowdisplaybreaks
\setlength{\emergencystretch}{2em}
\newtheorem{theorem}{Theorem}[section]
\newtheorem{lemma}[theorem]{Lemma}
\newtheorem{proposition}[theorem]{Proposition}
\newtheorem{corollary}[theorem]{Corollary}
\theoremstyle{remark}\newtheorem{remark}[theorem]{Remark}
\newcommand{\R}{\mathbb R}
\newcommand{\Sph}{\mathbb S^2}
\newcommand{\Vol}{\operatorname{Vol}}
\newcommand{\Q}{\mathcal Q}
\newcommand{\tr}{\operatorname{tr}}
\newcommand{\B}{\mathcal B}
\newcommand{\dd}{\,\mathrm d}
\newcommand{\VM}{V_{\mathrm M}}
\numberwithin{equation}{section}
\title{\Large\bfseries A rational quadratic support formulation for the sharp gap\\[5pt]
\normalsize Exact finite realizability, a fixed energy matrix, and an explicit volume error}
\author{}\date{23 September 2026}
\begin{document}\maketitle
\begin{abstract}
We replace the volume of a sample convex hull by a fixed quadratic objective
on finitely many support values. The support points satisfy exact antipodal
and diameter constraints, so every feasible finite array has a genuine
constant-width completion. On an antipodal spherical triangulation, the
piecewise linear homogeneous interpolant has a Dirichlet energy whose entire
matrix is rational when the normals are rational. A boundary integration
formula eliminates all angular integrals from that matrix. Uniform curvature
bounds give an explicit quadratic-in-mesh error for its volume objective.
The finite optimization is a concave quadratic minimization over a convex
second-order-cone feasible set; it is not a convex optimization problem.
Its exact global optima converge to the true constant-width minimum on
shape-regular refinements. We give a rational sufficient certificate for an
energy bound and exact checks of two finite matrices. No competitive global
minimum is computed here, and no new numerical lower bound or sharpness
claim is made.
\end{abstract}

\section{The objective: keep actual support data and remove hull combinatorics}
Let $K\subset\R^3$ have constant width one. Its support function satisfies
$h(n)+h(-n)=1$. Put
\[
 u=h-\frac12,\qquad u(-n)=-u(n),\qquad
 \Q u=\nabla^2_{\Sph}u+uI.
\]
The curvature-radius tensor of $K$ lies between zero and the identity:
\begin{equation}\label{eq:curvature}
 -\frac12 I\preceq\Q u\preceq\frac12 I.
\end{equation}
These statements have an almost-everywhere interpretation for $C^{1,1}$
support functions. The regularity follows, for example, from
$K+(-K)=\overline B(0,1)$: both support functions are convex summands of the
unit-ball support function, giving bounded second derivatives on spherical
charts. Standard support and mixed-volume facts are as in \cite{Schneider}.

Define the quadratic form
\begin{equation}\label{eq:energy}
 E(v)=\int_{\Sph}\bigl(|\nabla v|^2-2v^2\bigr)\dd\sigma,
 \qquad
 \mathcal V(v)=\frac\pi6-\frac14 E(v).
\end{equation}
For the actual centered support $u$,
\begin{equation}\label{eq:volume}
 \Vol(K)=\mathcal V(u).
\end{equation}
For completeness, the Codazzi property of $\Q h$ gives
$\operatorname{div}\operatorname{cof}(\Q h)=0$. Integration by parts then
shows
\[
 S(K)=\int\det(\Q h)
     =\int\left(h^2-\frac12|\nabla h|^2\right).
\]
Together with Blaschke's identity $\Vol(K)=S(K)/2-\pi/3$, and the vanishing
mean of the odd function $u$, this proves \cref{eq:volume}.
The eigenvalues of $-\Delta_{\Sph}$ are $l(l+1)$; hence $E(v)\ge0$ for every
odd $v\in H^1(\Sph)$, with equality precisely on linear functions
$v(n)=a\cdot n$. This also follows by applying the mean-zero spherical
Poincare inequality, whose equality space is the degree-one harmonics.

The proposal below has no independently assigned random curvature. Its
finite variables are actual support points. The only error is the explicit
interpolation error proved in \cref{sec:error}. The remaining difficulty is
a global finite quadratic optimization, not an unmeasured geometric
realizability error.

\section{Exact finite support interpolation}
Let $N=\{n_i\}\subset\Sph$ be finite and antipodally closed, with
$n_{\bar i}=-n_i$. Prescribe points $z_i\in\R^3$ and optionally a fixed finite
set $Y$ of diameter at most one.

\begin{lemma}[Diameter-one completion]\label{lem:completion}
Every nonempty finite set of diameter at most one is contained in a convex
body of constant width one.
\end{lemma}
\begin{proof}
For a current finite set $Z$, let
$C=\B(Z)=\bigcap_{z\in Z}\overline B(z,1)$.
The minimum containing ball of $Z$ has radius at most $\sqrt{3/8}$: its
center is a positive convex combination of at most four contact points,
and the weighted squared-distance identity gives
$R^2\le(1-\sum\lambda_i^2)/2\le3/8$.
Thus $C$ has a strictly feasible point.

For a unit $n$, choose $x$ maximizing $n\cdot x$ on $C$. The active-normal
rule gives $n=\sum a_i(x-z_i)$, where $a_i\ge0$, $|x-z_i|=1$, and
$\sum a_i\ge1$. If $w\in C$, then
\[
 (x-z_i)\cdot(w-x)\le-\frac12|w-x|^2,
\]
so $|w-(x-n)|^2\le1$. Consequently $Z\cup\{x,x-n\}$ still has diameter at
most one: $x$ lies in every old ball, $Z\subset C$, and $x-n\in\B(C)$.
Adjoin such a pair successively in a countable dense sequence of directions.
The sets stay in a fixed unit ball about an original point. The closed
convex hull of their union has diameter at most one and has width one in
every selected direction. Continuity of width in direction proves the claim.
\end{proof}

\begin{theorem}[Exact interpolation]\label{thm:interpolation}
A width-one body containing $Y$ has support point $z_i$ in direction $n_i$
for every $i$ if and only if
\begin{equation}\label{eq:finiteconstraints}
 \boxed{z_i-z_{\bar i}=n_i,\qquad
 |z_i-z_j|\le1,\qquad |z_i-y|\le1\quad(y\in Y),}
\end{equation}
with the already assumed diameter constraints within $Y$.
\end{theorem}
\begin{proof}
Opposite support points of a width-one body have scalar separation one in
their normal direction and Euclidean distance at most one. Their difference
is therefore exactly that unit normal, proving necessity.
Conversely apply \cref{lem:completion} to the finite set in
\cref{eq:finiteconstraints}. If $x$ belongs to the resulting body, then
\[
 n_i\cdot x\le n_i\cdot z_{\bar i}+|x-z_{\bar i}|
 \le n_i\cdot z_{\bar i}+1=n_i\cdot z_i.
\]
As $z_i$ is in the body, it is the prescribed support point.
\end{proof}

The feasible set is convex. Every constraint in \cref{eq:finiteconstraints}
is affine or a second-order-cone constraint. Its support values and centered
values are
\[
 h_i=n_i\cdot z_i,\qquad u_i=h_i-\frac12,
 \qquad u_{\bar i}=-u_i.
\]
For the unrestricted problem, include $\pm e_3$ in $N$ and impose
$z_{e_3}=e_3/2$, $z_{-e_3}=-e_3/2$. Every body admits this translation,
and it makes the finite feasible set compact. With a nonempty fixed seed,
compactness follows instead from a unit ball about one seed point.

\section{A fixed rational energy matrix}
Let $\mathcal T$ be an antipodally symmetric triangular surface with vertices
$n_i\in\Sph$, such that radial projection maps its faces bijectively, up to
shared edges, onto a triangulation of $\Sph$. For example, triangulate the
facets of an origin-containing inscribed convex polyhedron, consistently at
antipodes. Every face lies in an open hemisphere.

For an oriented face $\sigma=(n_0,n_1,n_2)$, order its vertices so that
$\det N_\sigma>0$, where $N_\sigma$ has rows $n_0^T,n_1^T,n_2^T$.
Let $\Omega_\sigma\subset\Sph$ be its radial image. Define
\[
 B_\sigma=N_\sigma^{-1}(u_0,u_1,u_2)^T,\qquad
 u_{\mathcal T}(n)=B_\sigma\cdot n\quad(n\in\Omega_\sigma).
\]
This interpolant is continuous: on a shared edge, the difference of the two
linear forms vanishes on the plane spanned by its two endpoints.
It is piecewise smooth, lies in $H^1(\Sph)$, and is odd.
It need not itself be a convex support function; exact realizability is
supplied by \cref{thm:interpolation}, not by imposing convexity on this
interpolant.

\begin{theorem}[Rational boundary formula]\label{thm:rationalmatrix}
For every face,
\begin{equation}\label{eq:faceenergy}
 \boxed{E_\sigma(u_{\mathcal T})
 =-\sum_{(i,j)=(0,1),(1,2),(2,0)}
 \frac{[B_\sigma\cdot(n_i\times n_j)]
       [B_\sigma\cdot(n_i+n_j)]}{1+n_i\cdot n_j}.}
\end{equation}
In particular, if all $n_i$ are rational, there is an explicitly computable
rational symmetric matrix $\mathsf E_{\mathcal T}$ such that
\begin{equation}\label{eq:matrixenergy}
 E(u_{\mathcal T})=\mathbf u^T\mathsf E_{\mathcal T}\mathbf u.
\end{equation}
On the antipodally odd nodal subspace this matrix is positive semidefinite.
Its kernel is exactly the three-dimensional space $u_i=a\cdot n_i$.
\end{theorem}
\begin{proof}
A linear spherical function $\ell(n)=B\cdot n$ satisfies $\Delta\ell=-2\ell$.
Green's identity on $\Omega_\sigma$ gives
\[
 \int_{\Omega_\sigma}(|\nabla\ell|^2-2\ell^2)
 =\int_{\partial\Omega_\sigma}\ell\,\partial_\nu\ell\,\dd s.
\]
The outward spherical conormal on the oriented great-circle edge from
$n_i$ to $n_j$ is the constant vector
$\nu=-(n_i\times n_j)/|n_i\times n_j|$.
If the edge length is $\theta\in(0,\pi)$, direct integration on its great
circle yields
\[
 \int_{\text{edge}}n\dd s=(n_i+n_j)\tan(\theta/2),\qquad
 \frac{\tan(\theta/2)}{|n_i\times n_j|}
 =\frac1{1+n_i\cdot n_j}.
\]
Also $\partial_\nu\ell=B\cdot\nu$. Substitution proves
\cref{eq:faceenergy}. Everything in this expression is rational in the
vertices and quadratic in their nodal values. Summing faces gives
\cref{eq:matrixenergy}. Positivity and the kernel follow from the spherical
Poincare inequality, since the assembled interpolant is odd. A zero-energy
interpolant is a global linear spherical function, and conversely every
such function is interpolated exactly.
\end{proof}

Here is an explicit matrix assembly formula. Put
\[
 T_\sigma=-\frac12\sum_{(i,j)}
 \frac{(n_i\times n_j)(n_i+n_j)^T+(n_i+n_j)(n_i\times n_j)^T}
 {1+n_i\cdot n_j}.
\]
The local matrix on the three nodal values is
$N_\sigma^{-T}T_\sigma N_\sigma^{-1}$.
Add it to the corresponding three-by-three block of the global matrix.
Only after assembly restrict to one value per antipodal pair.
A single local matrix need not be positive semidefinite.

The cancellation of all angular integrals is useful: the program's
quadratic coefficients are exact fractions, not numerical surface integrals.
In contrast, direct convex-hull volume changes its algebraic expression
when the sample hull changes facets. Here the triangulation and quadratic
matrix are fixed before the unknown support data are chosen.

\section{An explicit volume error independent of the unknown body}\label{sec:error}
For each flat face, define
\[
 c_\sigma=\operatorname{dist}(0,\operatorname{aff}\sigma)>0,\qquad
 d_\sigma=\max_{i,j\in\sigma}|n_i-n_j|,
\]
\begin{equation}\label{eq:errors}
 \epsilon_\sigma=\frac{1-c_\sigma^2}{4c_\sigma^2},\qquad
 b_\sigma=\frac{\sqrt3\,d_\sigma^2}
 {4c_\sigma\,\sigma_{\min}(N_\sigma)},\qquad
 \epsilon=\max_\sigma\epsilon_\sigma,\quad b=\max_\sigma b_\sigma.
\end{equation}
Here $\sigma_{\min}$ denotes the smallest singular value.
No unknown support function occurs in these constants.

\begin{lemma}[Uniform interpolation estimates]\label{lem:interperror}
For any actual width-one centered support $u$ interpolating the nodal data,
\begin{equation}\label{eq:interperror}
 \|u_{\mathcal T}-u\|_\infty\le\epsilon,
 \qquad |\nabla(u_{\mathcal T}-u)|\le b
 \quad\text{almost everywhere}.
\end{equation}
\end{lemma}
\begin{proof}
Extend $u$ homogeneously by $\widehat u(x)=|x|u(x/|x|)$.
Its Euclidean Hessian has radial eigenvalue zero and tangential block
$|x|^{-1}\Q u(x/|x|)$. By \cref{eq:curvature},
\[
 \|D^2\widehat u(x)\|\le\frac1{2|x|}.
\]
On a flat face, $|x|\ge c_\sigma$, and the same holds on every segment in
that face. Hence $\nabla\widehat u$ is Lipschitz there with constant
$1/(2c_\sigma)$.

Write $x=\sum\lambda_i n_i$, with $\lambda_i\ge0$, $\sum\lambda_i=1$, and
$r=|x|$. The affine interpolation remainder is bounded by
\[
 \left|\sum_i\lambda_i u_i-\widehat u(x)\right|
 \le\frac1{4c_\sigma}\sum_i\lambda_i|n_i-x|^2
 =\frac{1-r^2}{4c_\sigma}.
\]
Since $u_{\mathcal T}(x/r)=r^{-1}\sum\lambda_i u_i$, division by
$r\ge c_\sigma$ proves the value estimate.

For the gradient, let $a=\nabla\widehat u(x)$. Euler's identity gives
$a\cdot x=\widehat u(x)$. Taylor's formula at each vertex thus gives
\[
 (B_\sigma-a)\cdot n_i
 =u_i-\widehat u(x)-a\cdot(n_i-x),
 \qquad |(B_\sigma-a)\cdot n_i|
 \le\frac{d_\sigma^2}{4c_\sigma}.
\]
Inverting $N_\sigma$ yields $|B_\sigma-a|\le b_\sigma$.
Spherical gradients are the orthogonal projections of these homogeneous
Euclidean gradients, so their difference satisfies the same bound.
\end{proof}

A purely algebraic replacement for the singular-value bound is
\begin{equation}\label{eq:betaexplicit}
 b_\sigma^2\le
 \frac{9d_\sigma^6}{8c_\sigma^2(\det N_\sigma)^2}.
\end{equation}
Indeed $\sigma_1(N_\sigma)\le\sqrt3$. Subtracting the rank-one matrix with
all rows $n_0^T$ gives $\sigma_2(N_\sigma)\le\sqrt2d_\sigma$.
Their product and the determinant identity then give
$\sigma_{\min}\ge|\det N_\sigma|/(\sqrt6d_\sigma)$.
Both $c_\sigma^2$ and the right side of \cref{eq:betaexplicit} are rational
when the vertices are rational.

\begin{theorem}[Uniform quadratic-objective error]\label{thm:volumeerror}
Every genuine completion of an exact finite array satisfies
\begin{equation}\label{eq:errorbound}
 \boxed{\mathcal V(u_{\mathcal T})-2\pi\epsilon
 \le\Vol(K)\le
 \mathcal V(u_{\mathcal T})+2\pi\epsilon+\pi b^2.}
\end{equation}
\end{theorem}
\begin{proof}
Let $w=u_{\mathcal T}-u$, which is odd, and let $L=-\Delta-2$.
By \cref{eq:curvature}, $Lu=-\tr\Q u$ lies in $[-1,1]$ almost everywhere.
Weak integration by parts is legitimate for $u\in C^{1,1}$ and
$w\in H^1$. Expanding the energy gives the exact identity
\[
 \Vol(K)-\mathcal V(u_{\mathcal T})
 =\frac12\int (Lu)w+\frac14 E(w).
\]
The first term has absolute value at most $2\pi\epsilon$.
Oddness gives $E(w)\ge0$, while
$E(w)\le\int|\nabla w|^2\le4\pi b^2$.
This proves both bounds. Notice that the lower bound does not require
shape regularity or the gradient-error constant $b$.
\end{proof}

On a shape-regular sequence of meshes with maximum edge length $h\to0$,
$c_\sigma\to1$ uniformly, $1-c_\sigma^2=O(h^2)$, and
$|\det N_\sigma|$ is comparable to $d_\sigma^2$. Consequently
$\epsilon+b^2=O(h^2)$.
The needed geometric condition can be checked directly from
\cref{eq:errors,eq:betaexplicit}; there is no appeal to regularity of an
unknown minimizing body. Antipodally symmetric rational perturbations of
shape-regular refinements preserve these bounds if their displacement is
sufficiently smaller than the mesh size.

\section{The finite optimization and a rational certificate target}
Let $\mathcal Z_{\mathcal T}$ be the compact feasible set in
\cref{eq:finiteconstraints}, with the diametral normalization when there is
no fixed seed. Define
\[
 Q_{\mathcal T}=\max_{z\in\mathcal Z_{\mathcal T}}
          \mathbf u(z)^T\mathsf E_{\mathcal T}\mathbf u(z),\qquad
 q_{\mathcal T}=\frac\pi6-\frac14Q_{\mathcal T}.
\]
The finite objective is convex in the support values when written as the
energy to be maximized. Equivalently its volume form is a concave quadratic
to be minimized. Either optimization is generally nonconvex.

\begin{theorem}[A finite quadratic enclosure of the true minimum]
Let $V_*$ be the unrestricted width-one minimum, or the minimum in the
class containing the prescribed seed. Then
\begin{equation}\label{eq:minbound}
 \boxed{q_{\mathcal T}-2\pi\epsilon\le V_*
 \le q_{\mathcal T}+2\pi\epsilon+\pi b^2.}
\end{equation}
In particular, if $\epsilon+b^2\to0$, then $q_{\mathcal T}\to V_*$.
\end{theorem}
\begin{proof}
Every body produces a feasible finite array. Its interpolated quadratic
value is at least $q_{\mathcal T}$, so the lower bound follows from
\cref{thm:volumeerror}. Conversely, a minimizing finite array exists by
compactness and continuity. It has an actual completion by
\cref{thm:interpolation}. The upper bound in \cref{thm:volumeerror}, applied
to that completion, proves the other inequality.
\end{proof}

The finite values $q_{\mathcal T}$ are not claimed to be monotone under
refinement. Monotone lower and upper sequences can be formed by taking,
respectively, running maxima of the certified lower endpoints and running
minima of the certified upper endpoints in \cref{eq:minbound}.

For a practical sufficient certificate, eliminate all affine equations and
write the remaining support-point variables as a vector $x$. The diameter
constraints take the form
\[
 g_e(x)=1-|D_ex+d_e|^2\ge0,
 \qquad \mathbf u=Ux+v.
\]
For rational normals and a rational fixed seed, all these coefficients and
$\mathsf E_{\mathcal T}$ are rational. Any exact identity
\begin{equation}\label{eq:sdpcert}
 B-(Ux+v)^T\mathsf E_{\mathcal T}(Ux+v)
 =\begin{pmatrix}1\\x\end{pmatrix}^{\!T}
 Z\begin{pmatrix}1\\x\end{pmatrix}
 +\sum_e\lambda_e g_e(x),
 \qquad Z\succeq0,\quad\lambda_e\ge0,
\end{equation}
proves $Q_{\mathcal T}\le B$. It therefore implies
\[
 V_*\ge\frac\pi6-\frac B4-2\pi\epsilon.
\]
Coefficient equality and a rational positive-semidefinite factorization of
$Z$ are finite verification tasks. This is only a sufficient certificate;
we do not assert that constant multipliers $\lambda_e$ always attain the
exact finite maximum. Branch-and-bound or stronger polynomial certificates
may be necessary. An optimizer's feasible array supplies no upper bound on
$Q_{\mathcal T}$.

To reach the sharp Meissner value by this formulation, it would suffice to
prove, for meshes with $\epsilon\to0$,
\[
 Q_{\mathcal T}\le E_M+\delta_{\mathcal T},\qquad
 \delta_{\mathcal T}\to0,\qquad
 E_M=4\left(\frac\pi6-\VM\right)
 =\pi\sqrt3\arccos\frac13-2\pi.
\]
No such family of upper certificates is produced here. The new result makes
the required objective a fixed rational quadratic form while preserving
finite realizability; it does not solve its global maximization.

\section{Exact algebraic checks and numerical status}
The accompanying standard-library script constructs two rational normal
sets by stereographic grids of orders one and two, with duplicates removed.
It enumerates all supporting facets exactly, triangulates them antipodally,
and verifies oriented edge cancellation, vertex use, Euler's identity, and
positive radial orientation. For each triangulation it constructs the
matrix in \cref{thm:rationalmatrix} using fractions only, restricts to odd
values, and performs an exact semidefinite elimination. It also checks
that all three translation vectors lie in its kernel.

For a nonspherical feasible test body, take
\[
 h(n)=\frac12+\frac1{10}n_1n_2n_3.
\]
If $P(X)=X_1X_2X_3$, then
$\Q(P|_{\Sph})=D^2P|_{T\Sph}-2PI$.
On the sphere, $\|D^2P\|\le2$ and $|P|\le1/5$, so the perturbation's
curvature norm is at most $6/25<1/2$. The displayed $h$ is therefore a
width-one support function. Its support points at rational normals are
rational:
\[
 z(n)=\left(\frac12-\frac15P(n)\right)n+\frac1{10}\nabla P(n).
\]
The program checks all finite diameter inequalities and antipodal identities
for these data. Since $P|_{\Sph}$ has harmonic degree three and
$\int P^2\dd\sigma=4\pi/105$, its exact volume is
\[
 \frac\pi6-\frac{\pi}{1050}.
\]
This is a test body, not a minimum. Its computed finite energy is a check of
the rational matrix construction; the rigorous error estimates above are
proved analytically and can be coarse on these two small meshes.

The executed checks give:
\begin{center}
\begin{tabular}{rrrrrr}
\toprule
Grid & Vertices & Triangles & Odd variables & Energy rank & Nullity\\
\midrule
1&14&24&7&4&3\\
2&38&72&19&16&3\\
\bottomrule
\end{tabular}
\end{center}
The two exact test energies are $16/2025$ and $368/54675$.
They are values for the displayed test body, not maxima over feasible arrays.
All $91$ and $703$ pairwise diameter inequalities, respectively, pass.
The complete rational matrices and triangulations are supplied as JSON data.

Run \texttt{python check\_rational\_energy.py}. The script prints and writes
its exact counts, rational energies, kernel dimensions, and geometric error
constants. It does not solve a global quadratic optimization, certify a new
volume lower bound, or verify the sharp inequality. The universal value
$0.4144$ is established in the separate main manuscript and its fully rerun
certificate; the sharp gap is still open in this continuation.

\begin{thebibliography}{9}\small
\bibitem{Schneider} R. Schneider, \emph{Convex Bodies: The Brunn--Minkowski
Theory}, second expanded edition, Encyclopedia of Mathematics and its
Applications 151, Cambridge University Press.
\url{https://doi.org/10.1017/CBO9781139003858}.
\bibitem{Hynd} R. Hynd, \emph{On extreme constant width bodies in $\R^3$},
arXiv:2501.16940. Background on support-function constraints; the finite
interpolation and error estimates used here are proved in the text.
\end{thebibliography}
\end{document}
